\documentclass[
 paper=A4,pagesize=automedia,fontsize=11pt,
 BCOR=10mm,DIV=16,
 twoside,headinclude,footinclude=false,
 bibtotocnumbered,          % show bibl. refs in TOC
 liststotoc,                % show image list in TOC
 listsleft,
 pointlessnumbers,          % no point after heading index
 cleardoublepage=empty      % no page number on vacant pages
]{scrbook}

% pdfx wants to be first
\usepackage[a-1b,latxmp]{pdfx}

% encoding
\usepackage[utf8]{inputenc}
\usepackage[T1]{fontenc}
% language
\usepackage[english]{babel}
% fonts
\setkomafont{disposition}{\normalcolor\bfseries} % keine serifenlose Schrift
\usepackage{mathpazo}
% for being able to copy text out of the PDF
\usepackage{cmap}
% layout
\setlength\parindent{0em}
\usepackage{scrpage2} \pagestyle{scrheadings}
                      \clearscrheadfoot
                      \ihead{\headmark}\ohead{\pagemark}
                      \automark[section]{chapter}
                      \setheadsepline{0.5pt}
\usepackage{setspace} \onehalfspacing
\deffootnote{1em}{1em}{\textsuperscript{\thefootnotemark }}
% math mode formatting etc.
\usepackage{amsmath,amsfonts,amssymb,amsthm,sistyle,xcolor,delim}
\SIstyle{German}
\usepackage{graphicx,tikz,calc,tabularx}
\usetikzlibrary{backgrounds,calc}
% presentation of floats
\usepackage{flafter,afterpage}
\usepackage[section]{placeins}
\usepackage[margin=8mm,font=small,labelfont=bf,format=plain]{caption}
\usepackage[margin=8mm,font=small,labelfont=bf,format=hang]{subcaption}
% rendering of algorithms
\usepackage{algorithmic}
\usepackage[chapter]{algorithm}
\newcommand\algorithmheader[1][]{%
 \expandafter\def\expandafter\algorithmicrequire{\begin{minipage}{6em}\textbf{#1}~\end{minipage}}\REQUIRE%
}
\renewcommand\algorithmicrequire{\textbf{Input:}}
\renewcommand\algorithmicensure{\textbf{Output:}}
\newcommand\SUFFIXFOR[1]{\textbf{for every}~#1}
% put footnotes below floats
\usepackage[bottom]{footmisc}

\newcounter{theorem}
\theoremstyle{definition}
\newtheorem{theorem}[theorem]{Theorem}
\newtheorem{lemma}[theorem]{Lemma}
\newtheorem{definition*}[theorem]{Definition}

\newenvironment{definition}%
 {\begin{definition*}\pushQED{\qed}}%
 {\popQED\end{definition*}}%

\numberwithin{theorem}{chapter}
\numberwithin{equation}{chapter}
\numberwithin{figure}{chapter}
\numberwithin{table}{chapter}

% hyperref is loaded as the last package (as recommended by its manual)
\hypersetup{hyperindex,pdfborder={0 0 0},pdfa}

%%%%%%%%%%%%%%%%%%%%%%%%%%%%%%%%%%%%%%%%%%%%%%%%%%%%%%%%%%%%%%%%%%%%%%%%%%%%%%%%

\delimdef\kla#1{\dleft(#1\dright)}
\delimdef\brk#1{\dleft[#1\dright]}
\delimdef\brc#1{\dleft\{#1\dright\}}
\delimdef\abs#1{\dleft|#1\dright|}
\delimdef\floor#1{\dleft\lfloor#1\dright\rfloor}
\delimdef\ceil#1{\dleft\lceil#1\dright\rceil}
\delimdef\dangle#1{\dleft\langle#1\dright\rangle}
\delimdef\skal#1#2{\dleft\langle#1,#2\dright\rangle}
\delimdef\norm#1{\dleft\|#1\dright\|}

\usepackage{stmaryrd}
\delimdef\lang#1{\dleft\llbracket#1\dright\rrbracket}
\delimdef\stepmap#1{\llparenthesis#1\rrparenthesis}

\newcommand\ub{\mathcal B}
\newcommand\ud{\mathcal D}
\newcommand\ug{\mathcal G}
\newcommand\uh{\mathcal H}
\newcommand\ul{\mathcal L}
\newcommand\um{\mathcal M}
\newcommand\up{\mathcal P}
\newcommand\ur{\mathcal R}
\newcommand\zn{\mathbb N}
\newcommand\zr{\mathbb R}
\newcommand\zz{\mathbb Z}
\newcommand\argmax{\qopname\relax m{argmax}}
\newcommand\eps{\varepsilon}
\newcommand\mle{\operatorname{mle}}
\newcommand\cmle{\operatorname{cmle}}
\newcommand\pos{\operatorname{pos}}
\newcommand\osucc{\operatorname{succ}}
\newcommand\rk{\operatorname{rk}}
\newcommand\cnt{\operatorname{count}}
\newcommand\suff{\operatorname{suff}}
\newcommand\car{\operatorname{car}}
\newcommand\cdr{\operatorname{cdr}}

%%%%%%%%%%%%%%%%%%%%%%%%%%%%%%%%%%%%%%%%%%%%%%%%%%%%%%%%%%%%%%%%%%%%%%%%%%%%%%%%

\begin{document}

\frontmatter

\begin{titlepage}
 \begin{center}
  \vspace*{5em}

  \begin{singlespace}\bfseries\Huge
   Training of Hidden Markov models as an instance of the expectation maximization algorithm
  \end{singlespace}

  \vspace*{5em}

  \begin{singlespace}\large
   Bachelorarbeit \\ zur Erlangung des Hochschulgrades \\ Bachelor of Science
  \end{singlespace}\medskip

  \vspace*{4em}

  vorgelegt von \\
  {\large Dipl.-Phys.~Stefan Majewsky} \\
  geboren am 06.11.1989 in Schwerin

  \vspace*{4em}

  \begin{singlespace}\large
   Technische Universität Dresden \\
   Fakultät Informatik \\
   Institut für Theoretische Informatik \\
   Lehrstuhl für Grundlagen der Programmierung
  \end{singlespace}

  \vspace*{4em}

  \begin{singlespace}
   Betreuer: Dipl.-Inf.~Kilian~Gebhardt \\
   Verantwortlicher Hochschullehrer: Prof.~Dr.-Ing.~habil.~Heiko~Vogler
  \end{singlespace}

 \end{center}

 \vspace*{5em}
 \begin{singlespace}
  \hspace{0.18\linewidth}
  \begin{minipage}{0.1\linewidth}
   \includegraphics[width=\linewidth]{img/by-nd.pdf}
  \end{minipage}
  \hspace{0.02\linewidth}
  \begin{minipage}{0.6\linewidth}\footnotesize\flushleft
   This work is licensed under a \href{https://creativecommons.org/licenses/by-nd/4.0/}{Creative Commons Attribution-NoDerivatives 4.0 International License}.
  \end{minipage}\\[-1em]
  \begin{center}\footnotesize
   Source code at: \url{https://github.com/majewsky/bachelor-thesis/}
  \end{center}
 \end{singlespace}
\end{titlepage}

\cleardoublepage
\tableofcontents
\mainmatter

\chapter{Introduction}

In natural language processing, it is frequently necessary to judge the
correctness of sentences generated by some algorithm. For example, a
speech-to-text translator that only recognizes individual words might produce
the following two candidate sentences for the same input audio:
\begin{align*}
 \text{He ate soup.} &&
 \text{He aid soup.}
\end{align*}

Both sentences sound the same, but the second candidate sentence should be
discarded because it is syntactically wrong. As another example, a translation
algorithm that translates from another language to English might generate the
following two candidate sentences for some input:
\begin{align*}
 \text{This is a small red ball.} &&
 \text{This is a red small ball.}
\end{align*}

Both sentences are syntactically correct, but the first sentence should be
preferred because it conforms to the customary rules for adjective ordering in
the English language. Finally, consider a spelling correction process that is
applied to the following sentence:
\begin{align*}
 \text{The design \textbf{an} construction of the system will take more than a year. \cite{kukich1992}}
\end{align*}

Possible replacements for the syntactically wrong word ``an'' include ``a'' and
``and''. In all these situations, a \emph{language model} can be employed to
choose the best result from a set of candidates.
\cite{stolcke2002,youngetal2005} A language model assigns a probability $p_\omega(v)$
to each sentence $v\in V^*$ (with words from the set $V$), such that
correct sentences receive a higher probability than incorrect sentences, e.~g.,
\[
 1 > p_\omega\mbig\kla{\text{He ate soup.}} > p_\omega\mbig\kla{\text{He aid soup.}} > p_\omega\mbig\kla{\text{He He He soup.}} > 0.
\]

The probability distribution $p_\omega$ is determined by a \emph{model
parameter} $\omega\in\Omega$. The language model defines $\Omega$ and describes
how to obtain $p_\omega$ for any model parameter $\omega$. The model parameter
is chosen by a suitable \emph{training algorithm} using a training corpus $c$
containing known good sentences, such that $\omega$ maximizes the likelihood of
this corpus, i.~e.,
\[
 \prod_{v\in c} p_\omega(v).
\]

\pagebreak

Many training algorithms for language models are instances of the
\emph{expectation-ma\-xi\-mi\-za\-tion (EM) algorithm} \cite{demlairub77}. Chapter 2
will introduce a general framework for EM algorithms that first appeared in
\cite{bucstuvog15}.

\section{N-gram models}

\begin{figure}[b]
 \centering
 \begin{tikzpicture}
  \begin{scope}[every node/.style={rectangle,draw,minimum height=1.8em,minimum width=4em}]
   \node [minimum width=2em] (v0) at (-3,1) {\strut \#};
   \node (v11) at (0,0) {\strut He};
   \node (v12) at (0,2) {\strut She};
   \node (v21) at (4,0) {\strut ate};
   \node (v22) at (4,2) {\strut drank};
   \node (v31) at (7,0) {\strut soup};
   \node (v32) at (7,2) {\strut tea};
   \node [minimum width=2em] (v4) at (10,1) {\strut \#};
  \end{scope}
  \begin{scope}[every node/.style={circle,inner sep=0.2ex},thick,>=stealth]
   \draw[->] (v0.east) -- node[auto,swap] {$0.5$} (v11.west);
   \draw[->] (v0.east) -- node[auto] {$0.5$} (v12.west);
   \draw[->] (v11.east) -- node[auto,swap] {$0.5$} (v21.west);
   \draw[->] ($(v12.east)!0.5!(v12.south east)$) -- node[auto,pos=0.35,swap] {$0.5$} ($(v21.north west)!0.5!(v21.west)$);
   \draw[->] ($(v11.east)!0.5!(v11.north east)$) -- ($(v22.south west)!0.5!(v22.west)$);
   \draw[->] (v12.east) -- node[auto] {$0.5$} (v22.west);
   \draw[->] (v21.east) -- node[auto,swap] {~$1$~} (v31.west);
   \draw[->] ($(v22.east)!0.5!(v22.south east)$) -- node[auto,swap] {$0.3$} ($(v31.north west)!0.5!(v31.west)$);
   \draw[->] (v22.east) -- node[auto] {$0.7$} (v32.west);
   \draw[->] (v31.east) -- node[auto,swap] {~$1$~} ($(v4.south west)!0.8!(v4.west)$);
   \draw[->] (v32.east) -- node[auto] {~$1$~} ($(v4.north west)!0.8!(v4.west)$);
  \end{scope}
 \end{tikzpicture}
 \caption{Graphic representation of an example instance of the bigram model.
 The nodes are words from the vocabulary $V\cup\{\#\}$. The start/end marker
 $\#$ is shown twice for readability's sake. Arrows $v\to v'$ represent a
 nonzero bigram probability $b(v'|v)$.\label{fig:01-bigram}}
\end{figure}

One of the simplest language models is the \emph{bigram model}.
\cite{jm09} It interprets the generation of a sentence as a stochastic
process, wherein each word is chosen with a probability conditional on the word
that appeared before it: (see~Figure~\ref{fig:01-bigram})
\[
 p(\text{He ate soup.}) = b(\text{He}|\#) \cdot b(\text{ate}|\text{He}) \cdot b(\text{soup}|\text{ate}) \cdot b(\#|\text{soup}).
\]
The model parameter is $b$, a conditional probability distribution of
$V\cup\brc{\#}$ given $V\cup\brc{\#}$, where $\#$ is a placeholder word
that stands in for the start and end of the sentence.

The model parameter can be chosen by a very simple training algorithm
\cite[pp.~123]{jm09}: All bigrams, i.~e., pairs of words occurring one after the
other, are counted across the corpus. This includes bigrams containing the
sentence delimiter word $\#$. Then, the counts are normalized into a
conditional probability distribution by setting
\[
 b(w|w') := \frac{\operatorname{count}(w'w)}{\sum_{w''\in V\cup\brc{\#}} \operatorname{count}(w'w'')}
 \quad\forall w,w'\in V\cup\brc{\#}.
\]

The bigram model can be generalized to the \emph{N-gram model}, wherein the
next word is chosen with a probability conditional on the $n-1$ words before
it. Training then counts $n$-grams, i.~e., sequences of $n$ words, hence the
name \emph{N-gram model}. The bigram model is recovered for $n=2$.

Besides the already mentioned applications where language models augment
translators and autocorrect systems, N-gram models are useful in
\emph{augmentive communication}: Virtual keyboards on smartphones and tablets
predict the next word by looking at the previous $n$ words, thus improving
typing speed and accuracy.  \cite{hasan2004n} A similar assistive word
prediction system is part of speech synthesis programs used by disabled persons
such as the physicist Stephen Hawking. \cite{newelletal1998}

\section{Hidden Markov model}

\begin{figure}[b!]
 \centering
 \begin{tikzpicture}[baseline=(current bounding box.north)]
  \begin{scope}[every node/.style={circle,draw}]
   \node[double] (q0) at (-3,3) {\#};
   \node (q1) at (0,3) {noun};
   \node (q2) at (0,0) {verb};
  \end{scope}
  \begin{scope}[every node/.style={anchor=west}]
   \node (v11) at (4,4.5) {he};
   \node (v12) at (4,3.5) {she};
   \node (v13) at (4,2.5) {soup};
   \node (v14) at (4,1.5) {tea};
   \node (v21) at (4,0.5) {drank};
   \node (v22) at (4,-0.5) {ate};
  \end{scope}
  \begin{scope}[every node/.style={rectangle,inner sep=0.5ex},thick,>=stealth]
   \draw[<-] (q0.west) -- +(-1,0);
   \draw[->] (q0) edge[bend left=20] node[auto,pos=0.55] {$1$} (q1);
   \draw[->] (q1) edge[bend right=20] node[auto,swap] {$0.5$} (q2);
   \draw[->] (q1) edge[bend left=20] node[auto,pos=0.45] {$0.5$} (q0);
   \draw[->] (q2) edge[bend right=20] node[auto,swap] {$1$} (q1);
   \draw[dashed,->] (q1) edge[bend left=10] node[auto,pos=0.63] {$0.3$} (v11);
   \draw[dashed,->] (q1) edge node[auto,pos=0.6] {$0.3$} (v12);
   \draw[dashed,->] (q1) edge node[auto,pos=0.63,swap] {$0.25$} (v13);
   \draw[dashed,->] (q1) edge[bend right=10] node[auto,pos=0.67,swap] {$0.15$} (v14);
   \draw[dashed,->] (q2) edge node[auto,pos=0.6] {$0.5$} (v21);
   \draw[dashed,->] (q2) edge node[auto,pos=0.6,swap] {$0.5$} (v22);
  \end{scope}
 \end{tikzpicture}
 \caption{Graphic representation of an example instance of the Hidden Markov
 model, with hidden states ``noun'' and ``verb''; adapted from
 \cite{nel13}. Arrows $q\to q'$ and dashed arrows $q\dashrightarrow v$
 represent nonzero transition and emission probabilities $t(q'|q)$ and
 $e(v|q)$, respectively.\label{fig:01-hmm}}
\end{figure}

A further generalization of N-gram models leads to the \emph{Hidden Markov
model}. In this model, the emission probability of a word depends not on the
previous words, but on the progression of a state machine that cannot be
directly deduced from the emitted words. Every time a word needs to be
emitted, the state machine progresses to a new state according to a transition
probability distribution $t$ dependent on the previous state, and the next word
is predicted by an emission probability distribution $e$ dependent on the new
state. The symbol $\#$ is used as the initial and final state of the state
machine.

For example, if the state sequence that generated the sentence ``He ate soup''
was ``\#-noun-verb-noun-\#'', then the probability of that sentence would be
\begin{align*}
 p(v|q)
  =&\; t(\text{noun}|\#) \cdot e(\text{he}|\text{noun}) \cdot t(\text{verb}|\text{noun}) \cdot e(\text{ate}|\text{verb}) \\
  &\cdot t(\text{noun}|\text{verb}) \cdot e(\text{soup}|\text{noun}) \cdot t(\#|\text{noun}).
\end{align*}

Since the state sequence is typically not known, a sum over all possible state
sequences $q$ must be computed to obtain $p(v)$ for a sentence $v\in V^*$.

N-gram models can be interpreted a special case of the Hidden Markov model by
encoding the current N-gram in the state ($Q = V^n$) and defining the
transmission and emission probability in terms of the N-gram probability. For
example, for bigrams,
\begin{align*}
 t(v_1v_2|v_1'v_2') &:= \begin{cases}
  b(v_2|v_1) & \text{if } v_1 = v_2', \\
  0 &\text{otherwise},
 \end{cases} &
 e(v|v_1v_2) &:= \begin{cases}
  1 & \text{if } v = v_2, \\
  0 &\text{otherwise}.
 \end{cases}
\end{align*}

Training for the Hidden Markov model is not as straightforward as for the
bigram model, since the training corpus typically only contains the observed
sentences, not the state sequences that produced them. The standard training
algorithm for Hidden Markov models is the \emph{Baum-Welch algorithm}.
\cite{baupetsouwei70,baum1972}

Chapter~3 will define Hidden Markov models more formally, and discuss
algorithms that act on them. Finally, we will show that the Baum-Welch
algorithm is an instance of the general EM algorithm laid out in Chapter~2 and
thus exhibits the same beneficial convergence properties.

\chapter{Expectation-maximization algorithms}

The generic framework for expectation-maximization (EM) algorithms introduced
by \cite{bucstuvog15} applies to both language models and translation models.
This chapter will introduce terminology and notation from \cite{bucstuvog15} as
far as is necessary to apply this framework to the Hidden Markov Model in the
subsequent chapters.

\section{Preliminaries}

The set $\brc{0,1,2,\ldots}$ of non-negative integers and the set of
non-negative reals shall be denoted by $\zn$ and $\zr_{\geq0}$, respectively.
We assume that
\begin{align*}
 0^0 &:= 1, &
 \log 0 &:= -\infty, &
 0 (-\infty) &= \log 0^0 = 1.
\end{align*}

\begin{definition}
 Given a countable set $X$, a mapping $c: X \to \zr_{\geq0}$ is called a
 \emph{$X$-corpus}.
\end{definition}

When used as input for a language model's training algorithm, $X$ is the set of
all sentences consisting of words from the language in question, and $c(x)$
describes how often a sentence $x\in X$ occurs in the corpus $c$. The
\emph{size} of the corpus is defined as
\begin{equation*}
 \abs c := \sum_{x\in X} c(x).
\end{equation*}

\begin{definition}
 A \emph{probability distribution of $X$} is an $X$-corpus of size 1.
\end{definition}

The set of all such probability distributions is denoted by $\um(X)$.
Probability distributions can be derived from corpora:

\begin{definition}
 Given a non-empty and finite $X$-corpus $c$, the \emph{empirical probability
 distribution} $\tilde c$ is defined as
 \[
  \tilde c(x) = \frac{c(x)}{\abs c}.
 \]
\end{definition}

This expression is not well-defined for $\abs c = 0$ or $\abs c = \infty$,
hence the requirement for $c$ to be non-empty and finite. If this cannot be
guaranteed, a fall-back function can be added.

\begin{definition}
 Given $p\in\um(X)$, a \emph{normalization mapping with fall-back $p$} is the
 mapping
 \[
  \overline p: \zr_{\geq0}^X \to \um(X),
  \quad
  c \mapsto \begin{cases}
   \tilde c & \text{if } 0 < \abs c < \infty, \\ p & \text{otherwise}.
  \end{cases}
 \]
\end{definition}

\begin{definition}
 Given an $X$-corpus $c$ and $p\in\um(X)$, the \emph{likelihood of $c$ under
 $p$} is
 \[
  p(c) := \prod_{x\in X} p(x)^{c(x)}.
 \]
\end{definition}

The likelihood describes the probability of observing the sentences from the
corpus $c$ when sentences occur with the probability distribution described by
$p$. A training algorithm will take $c$ as an input, and seek to find an
admissible $p$ such that $p(c)$ is maximized. If any $p$ is admissible, then
for non-empty and finite $c$, the optimal choice is $p = \tilde c$:

\begin{lemma}\label{lemma:empirical1}
 Let $c$ be a non-empty and finite $X$-corpus. Then $\tilde c(c) \geq p(c)$ for
 every $p\in\um(X)$.
\end{lemma}

\begin{proof}
 Since $\log$ is monotone, it suffices to show that $\log\tilde c(c) \geq \log
 p(c)$. Using Gibbs' inequality,
 \begin{align*}
  \log \tilde c(c)
  &= \sum_{x\in X} c(x) \cdot \log \tilde c(x)
  = \abs c \cdot \sum_{x\in X} \tilde c(x) \cdot \log \tilde c(x) \\
  &\geq \abs c \cdot \sum_{x\in X} \tilde c(x) \cdot \log p(x)
  = \sum_{x\in X} c(x) \cdot \log p(x)
  = p(c)
  \qedhere
 \end{align*}
\end{proof}

However, using $p = \tilde c$ directly is not useful because this probability
distribution is grossly overfitted: It will assign zero probability to any
sentence not in the original corpus. A useful language model thus limits the
set of admissible $p$ by describing the probability distribution in terms of
\emph{model parameters} $\omega\in\Omega$.

\begin{definition}
 Given a set $\Omega$, an \emph{$\Omega$-probability model for $X$} is a
 mapping $p:\Omega\to\um(X)$.
\end{definition}

Instead of $p(\omega)$, we write $p_\omega$. Training shall then find
$\omega\in\Omega$ such that $p_\omega(c)$ is maximized.

\begin{definition}
 Given a set $\Omega$ and an $\Omega$-probability model $p$ for $X$, the
 \emph{maximum likelihood estimator} for $p$ is the mapping
 \[
  \mle_p: \zr_{\geq0}^X \to \up(\Omega),
  \quad
  c \mapsto \argmax_\omega p_\omega(c).
 \]
\end{definition}

$\mle_p(c)$ is the set of all $\omega$ with maximal likelihood, but training
only needs to find a single $\hat\omega \in \mle_p(c)$. Computing $\mle_p(c)$
by brute force is typically not tractible because the set $\Omega$ is infinite
(usually countably infinite). However, there is one easily solvable special
case.

\begin{lemma}\label{lemma:empirical2}
 Let $c$ be a finite $X$-corpus, and $p$ a $\Omega$-probability model for $X$.
 If there exists $\hat\omega\in\Omega$ such that $p_{\hat\omega} = \tilde c$,
 then $\hat\omega\in\mle_p(c)$.
\end{lemma}

\begin{proof}
 If $c$ is empty, then $p_\omega(c) = 1$ for every $\omega$, and thus
 $\mle_p(c) = \Omega \ni \hat\omega$. Otherwise, by
 Lemma~\ref{lemma:empirical1}, $p_{\hat\omega}(c) = \tilde c(c) \geq
 p_\omega(c)$ for every $\omega\in\Omega$, and thus $\hat\omega \in \mle_p(c)$.
\end{proof}

\section{Algorithmic skeleton}

\begin{algorithm}[t]
 \caption{Algorithmic skeleton for EM of language models according to \cite{bucstuvog15}}
 \label{alg:skeleton}
 \begin{algorithmic}[1]
  \algorithmheader[Input:] $X$-corpus $c$
  \algorithmheader         $\Omega$-probability model $p$ for $Y\times X$
  \algorithmheader         some initial parameter $\omega_0 \in \Omega_0$ where $\Omega_0 := \brc{\omega\in\Omega: p_\omega(c) \neq 0}$
  \algorithmheader[Implicit:] step mapping $\psi:\Omega_0\to\up(\Omega)$
  \algorithmheader            \hspace{1em} such that $\omega'\in\psi(\omega)$ implies $p_\omega(c) \leq p_{\omega'}(c)$
  \algorithmheader[Output:] sequence $\omega_1,\omega_2,\ldots\in\Omega_0$
  \algorithmheader            \hspace{1em} such that $p_{\omega_0}(c), p_{\omega_1}(c), p_{\omega_2}(c)$ nondecreasing

  \STATE $i\leftarrow 0$
  \WHILE{not converged}
   \STATE $\omega_{i+1} \leftarrow \text{select a member of $\psi(\omega_i)$}$
   \STATE output $\omega_{i+1}$
   \STATE $i\leftarrow i+1$
  \ENDWHILE
 \end{algorithmic}
\end{algorithm}

The major complication that occurs when trying to train a typical language
model is that not all required data is present in the corpus. For example, a
probabilistic context-free grammar is described by the probability distribution
of derivation rules. \cite{laryou90} When the training data consists of full
parse trees (\emph{supervized training}), the optimal probability distribution
can be found by simply counting how many times each rule is used across all
these parse trees, and then computing the empirical probability distribution
for this corpus. Most of the times, however, the training data will consist
only of sentences (\emph{unsupervized training}). The information about how to
parse the sentences is hidden.

The same problem arises with the Hidden Markov Model: When training data is not
already annotated with state information, the information which states
correspond to which words from the training data remains hidden. Expectation
maximization algorithms can be used when parts of the training data are hidden
in such a way.

For the remainder, let
\begin{itemize}\setlength\itemsep{-0.3em}
 \item $X$ and $Y$ be countable sets,
 \item $\Omega$ be a set,
 \item $c$ be a finite $X$-corpus and
 \item $p$ be a $\Omega$-probability model for $Y\times X$.
\end{itemize}

$c$ represents the set of training data. Each $x\in\operatorname{supp}(c)$ is
an \emph{observation}. To judge its probability under a $p_\omega$, additional
\emph{hidden information} $y\in Y$ is required. For notational convenience, we define
\begin{align*}
 p_\omega(c) &:= \prod_{x\in X} p_\omega(x)^{c(x)}, &
 \text{where } p_\omega(x) &:= \sum_{y\in Y} p_\omega(x,y).
\end{align*}

That is, even though $p_\omega$ is a probability distribution over $Y\times X$,
we allow to take the likelihood of the $X$-corpus $c$ under $p$ by aggregating
the probabilities for all hidden information $y$ that lead to a certain
observation $x$.

The basic pattern for expectation maximization is outlined in
algorithm~\ref{alg:skeleton}. The algorithm starts with an initial $\omega_0$
such that $p_{\omega_0}(c) \neq 0$. It then iteratively employs a step mapping
to choose the next $\omega_i$ with a higher (or at least equal) likelihood than
the one that came before.

\begin{definition}
 A \emph{step mapping} is a mapping $\psi:\Omega_0\to\up(\Omega)$ which is
 nondecreasing in the following manner:
 \[
  \forall \omega\in\Omega_0: \forall \omega'\in\psi(\Omega): p_\omega(c) \leq p_{\omega'}(c)
 \]
\end{definition}

The step mappings that we will consider will typically consist of two steps:
\begin{enumerate}
 \item \emph{Expectation:} The training data $c$ is converted into a
  \emph{complete-data corpus}. Using the $\omega_i$ from the previous
  iteration, the complete-data corpus estimates how hidden information
  contributes to the observations in the original corpus.
 \item \emph{Maximization:} A suitable maximum-likelihood estimator is applied
  to the complete-data corpus to choose $\omega_{i+1}$.
\end{enumerate}

%TODO: this section may use some citations because it makes claims
This back and forth of using the current $\omega$ to enrich the training data
and using the enriched data to find a better $\omega$ will converge towards a
local maximum of likelihood. The iteration is therefore usually aborted after
the desired running time has been exceeded, or after the changes of
$p_{\omega_i}(c)$ per iteration have become smaller than some threshold.

\cite{bucstuvog15} identify three types of step mappings that build on each
other, each one more specific than the one before it. Since the training of
Hidden Markov Models will be identified as an instance of the most specific
step mapping, the remainder of this chapter will introduce all three in order.

\section{Corpus-based step mapping}

The most general type of complete-data corpus can be obtained by distributing
$c(x)$ among the hidden information $y$ according to the probability
distribution $p_\omega$:
\[
 c\!\dangle{\omega,p}(y,x) := \begin{cases}
  c(x) \cdot \frac{p_\omega(y,x)}{p_\omega(x)} & \text{if } p_\omega(x) \neq 0, \\
  0 & \text{if } p_\omega(x) = 0.
 \end{cases}
\]
Recall that $p_\omega(x) = \sum_y p_\omega(y,x)$. Therefore,
$\abs{c\!\dangle{\omega,p}} = \abs c$. This corpus now has the correct
structure for plugging it into $\mle_p$, yielding the \emph{corpus-based step
mapping}\footnote{The proof that this step mapping is nondecreasing can be
found in \cite[pp.~10]{bucstuvog15}.}
\[
 \stepmap p_\mathrm{cb}: \Omega_0\to\up(\Omega),
 \quad
 \omega \mapsto \mle_p\mbig\kla{c\!\dangle{\omega,p}} = \argmax_{\omega'} p_{\omega'}\mbig\kla{c\!\dangle{\omega,p}}.
\]

For very simple language models, the $\argmax$ can be solved at this point
already. To apply lemma~\ref{lemma:empirical2}, $\hat\omega$ needs to be found
such that $p_{\hat\omega} = \widetilde{c\!\dangle{\omega,p}}$. This operation
is typically hard, which is why a more specific step mapping is helpful.

\section{Simple counting step mapping}

The next such step mapping requires the language model to be described by a
\emph{counting information}. Before defining this term, some additional
notation needs to be introduced.

\begin{definition}
 Let $A$ and $B$ be sets. A mapping $p: B \to \um(A)$ is called
 \emph{conditional probability distributions of $A$ given $B$}. The set of all
 such mappings is denoted by $\um(A|B)$.
\end{definition}

To simply notation, we define $p(a|b) := p(b)(a)$.

\begin{definition}
 Let $C\subseteq A\times B$ be a set. Then
 \[
  \um_C(A|B) := \brc{p\in\um(A|B): \operatorname{supp}(p)\subseteq C}
 \]
 is the set of all conditional probability distributions of $A$ given $B$
 constrained to $C$.
\end{definition}

\begin{definition}
 Given a set $\Omega$, a \emph{conditional $\Omega$-probability model for $A$
 and $B$ (constrained to $C$)} is a mapping $q:\Omega\to\um(A|B)$ (or
 $q:\Omega\to\um_C(A|B)$).
\end{definition}

When talking about counting informations (and inside-outside informations in
the next section), we assume the previously established requirements for $X$,
$Y$, $\Omega$, $c$ and $p$. Furthermore, we require that $X$ and $Y$ both
contain a special symbol $\bot$ such that $c(\bot) = 0$. $\bot$ can easily be
added to any previously defined $X$ and $Y$ without affecting the requirement
for countability. The notation $U_{\not\bot} := U\setminus\brc\bot$ shall be
defined for any set $U$.

\begin{definition}
 Let $A$, $B$ and $C$ be sets such that $C\subseteq A\times B$. A
 \emph{counting information} is a triple $\varkappa = (q,\lambda,\pi)$ such
 that
 \begin{align*}
  q &: \Omega\to\um_C(A|B), &
  \lambda &: X_{\not\bot} \times Y_{\not\bot} \to [0,1], &
  \pi &: X_{\not\bot} \times Y_{\not\bot} \to \zr_{\geq0}^C.
 \end{align*}
\end{definition}

The motivation for this definition is to model hidden information $y$ as
consisting of countable events $c\in C$. $\lambda(x,y)$ describes whether (and,
possibly, with what probability) a certain $y$ can be the cause for a certain
observation $x$.\footnote{Most actual instances of $\lambda$ use only integer
images, i.~e.~$\lambda(X_{\not\bot}\times Y_{\not\bot}) = \brc{0,1}$, thus
following this intuitive notion. However, the possibility of using fractional
values for $\lambda(x,y)$ is occasionally useful, e.~g.~to define a counting
information for the IBM Model 1 in \cite[pp.~23]{bucstuvog15}.} For
$\lambda(x,y)>0$, $\pi(x,y)$ is a $C$-corpus that describes how often each
countable event occurs in this hidden information.

Following the assumption that the countable events $C$ fully encode the hidden
information $Y$, we can use these intuitive notions to describe the original
probability model $p$ in terms of the counting information.

\begin{definition}
 Given $\varkappa=(q,\lambda,\pi)$, the \emph{induced model}
 $\varkappa^\flat:\Omega\to\zr_{\geq0}^{Y\times X}$ is given by
 \[
  (\varkappa^\flat)_\omega(y,x) := \begin{cases}
   \lambda(x,y) \cdot q_\omega\mbig\kla{\pi(x,y)} & \text{if } x,y\neq\bot, \\
   1 - \sum_{x',y'\neq\bot} \lambda(x',y') \cdot q_\omega\mbig\kla{\pi(x',y')} & \text{if } x = y = \bot, \\
   0 & \text{otherwise}.
  \end{cases}
 \]
\end{definition}

This definition shows why the introduction of $\bot$ into $X$ and $Y$ was
useful. By defining $(\varkappa^\flat)_\omega(\bot,\bot)$ as above, we ensure
$\mnorm\abs{(\varkappa^\flat)_\omega} = 1$. Therefore, $\varkappa^\flat$ is an
$\Omega$-probability model for $Y\times X$ iff
$(\varkappa^\flat)_\omega(\bot,\bot) \geq 0$. We call $\varkappa$ \emph{proper}
in this case.

Since we now have a probability model $p = \varkappa^\flat$ as required by the
corpus-based step mapping, we can lift its complete-data corpus into the domain
of the counting information, obtaining a new complete-data corpus
\[
 c\!\dangle{\omega,\varkappa}: C \to \zr_{\geq0}, \quad
 (a,b) \mapsto \sum_{x,y} c\mnorm\dangle{\omega,\varkappa^\flat}(y,x) \cdot \pi(x,y)(a,b).
\]

We can now apply a maximum-likelihood estimator for $q$ to arrive at the step
mapping for this class of language models.

\begin{definition}
 Given a set $\Omega$ and a conditional $\Omega$-probability model $q$ for $A$
 and $B$, the \emph{conditional maximum likelihood estimator} for $q$ is the mapping
 \[
  \cmle_q: \zr_{\geq0}^{A\times B} \to \up(\Omega),
  \quad
  c \mapsto \argmax_\omega q_\omega(c).
 \]
\end{definition}

Using this definition, the \emph{simple counting step mapping}\footnote{The
proof that this step mapping is equivalent to
$\mnorm\stepmap{\varkappa^\flat}_\mathrm{cb}$, and thus also nondecreasing, can
be found in \cite[p.~13]{bucstuvog15}.} is
\[
 \stepmap\varkappa_\mathrm{sc}: \Omega_0 \to \up(\Omega), \quad
 \omega \mapsto \cmle_q\mbig\kla{c\mnorm\dangle{\omega,\varkappa}} = \argmax_{\omega'} q_{\omega'}\mbig\kla{c\mnorm\dangle{\omega,\varkappa}}.
\]
The simple counting step mapping has two advantages over
$\stepmap\cdot_\mathrm{cb}$: First, many language models can be described in
terms of the countable events only, such that $\Omega = C$. In this case,
lemma~\ref{lemma:empirical2} can be used to solve the $\argmax$ by simply
computing the empirical probability distribution of
$c\!\dangle{\omega,\varkappa}$.

Second, even if this is not possible, the set of countable events $C$ is
usually much smaller than the set of all observations $X$ or hidden information
$Y$, making the evaluation of the $\argmax$ more tractable than for the
corpus-based step mapping. For example, when considering probabilistic
context-free grammars, $X$ (the set of all sentences) and $Y$ (the set of all
parse trees) are both countably infinite, but $C$ (the set of all derivation
rules) is finite.

\section{Regular tree grammars}

The third and most specific type of step mapping applies to language models
whose hidden information can be described as trees, such that the countable
events are labels in these trees' nodes. We therefore need to define a set of
terms regarding trees and tree grammars first.

\begin{definition}
 An \emph{alphabet} is a finite set.
\end{definition}

\begin{definition}
 Let $\Sigma$ be an alphabet and $V$ be a set. The \emph{set $U_\Sigma(V)$ of
 unranked trees over $\Sigma$ indexed by $V$} is the smallest set $T$ such that
 \[
  V \subseteq T \quad\wedge\quad \forall k\in\zn: \forall \sigma\in\Sigma, t_1,\ldots,t_k\in T: \sigma(t_1,\ldots,t_k)\in T.
 \]
\end{definition}

The notation $\sigma(t_1,\ldots,t_k)$ refers to the tree which has the label
$\sigma$ at its root and the subtrees $t_1,\ldots,t_k$ in that order. For
$k=0$, we shorten $\sigma() \equiv \sigma$. For $V=\emptyset$, we shorten
$U_\Sigma(\emptyset) \equiv U_\Sigma$. Trees from $U_\Sigma(V)$ have labels
from $\Sigma\cup V$ at each node, but labels from $V$ are only permitted at
leafs. For example, for $\Sigma = \brc{\sigma_1,\sigma_2}$ and
$V=\brc{v_1,v_2}$:

\[\begin{matrix}
 \begin{tikzpicture}[every node/.style={rectangle,draw,minimum size=2em},baseline=(current bounding box.center)]
  \node (n)   at (+0,+0.0) { $\sigma_1$ };
  \node (n1)  at (-1,-1.5) { $\sigma_2$ } edge (n);
  \node (n2)  at (+1,-1.5) { $\sigma_2$ } edge (n);
  \node (n11) at (-2,-3.0) { $v_1$ } edge (n1);
  \node (n12) at (+0,-3.0) { $\sigma_1$ } edge (n1);
  \node (n21) at (+2,-3.0) { $v_2$ } edge (n2);
 \end{tikzpicture}
 &\hspace*{4em}&
 \begin{tikzpicture}[every node/.style={rectangle,draw,minimum size=2em},baseline=(current bounding box.center)]
  \node (n)   at (+0,+0.0) { $\sigma_1$ };
  \node (n1)  at (-1,-1.5) { $\sigma_2$ } edge (n);
  \node (n2)  at (+1,-1.5) { $v_1$ } edge (n);
  \node (n11) at (-2,-3.0) { $v_1$ } edge (n1);
  \node (n12) at (+0,-3.0) { $\sigma_1$ } edge (n1);
  \node (n21) at (+2,-3.0) { $v_2$ } edge (n2);
  \draw (n2) circle (0.8);
  \node[draw=none] at (+1.75,-0.75) {!};
 \end{tikzpicture}
 \\\\
 \sigma_1\mbig\kla{\sigma_2(v_1,\sigma_1), \sigma_2(v_2)} \in U_\Sigma(V) &&
 \sigma_1\mbig\kla{\sigma_2(v_1,\sigma_1), v_1(v_2)} \notin U_\Sigma(V)
\end{matrix}\]

We refer to positions in the tree using \emph{Gorn notation}. As a quick
illustration, in the figure below, each node of the left tree is labeled with
its position in the tree:

\vspace*{-1em}\[
 \hspace*{3em}
 \begin{tikzpicture}[every node/.style={rectangle,draw,minimum size=2em},baseline=(current bounding box.center)]
  \node (n)   at (+0,+0.0) { $\eps$ };
  \node (n1)  at (-1,-1.5) { $1$ } edge (n);
  \node (n2)  at (+1,-1.5) { $2$ } edge (n);
  \node (n11) at (-2,-3.0) { $1\ 1$ } edge (n1);
  \node (n12) at (+0,-3.0) { $1\ 2$ } edge (n1);
  \node (n21) at (+2,-3.0) { $2\ 1$ } edge (n2);
 \end{tikzpicture}
 \hspace*{6em}
 \begin{tikzpicture}[baseline=(current bounding box.center)]
  \begin{scope}[every node/.style={rectangle,draw,fill=white,minimum size=2em}]
   \node (n)   at (+0,+0.0) { $\sigma_1$ };
   \node (n1)  at (-1,-1.5) { $\sigma_2$ } edge (n);
   \node (n2)  at (+1,-1.5) { $\sigma_2$ } edge (n);
   \node (n11) at (-2,-3.0) { $v_1$ } edge (n1);
   \node (n12) at (+0,-3.0) { $\sigma_1$ } edge (n1);
   \node (n21) at (+2,-3.0) {} edge (n2);
  \end{scope}
  \begin{scope}[every node/.style={inner sep=0},every pin/.style={inner sep=0},pin distance=1cm]
   \node[draw=none,pin=60:{$t(2\ 1)$}] at (n21.center) { $v_2$ };
  \end{scope}
  \begin{scope}[on background layer]
   \draw[line join=round,draw=black!10,fill=black!10,line width=1.5cm] (n1.center) -- (n11.center) -- (n12.center) -- cycle;
  \end{scope}
  \node at (-2.4,-1.7) { $t|_1$ };
 \end{tikzpicture}
\]

$\eps$ is the empty word. The set of all positions in a tree $t$ shall be
denoted by $\pos(t)$. The right tree illustrates some notation that we're going
to use for trees. For any trees $t,t'$ and position $w\in\pos(t)$,
\begin{itemize}\setlength\itemsep{-0.3em}
 \item $t(w)$ is the label at $w$,
 \item $t|_w$ is the subtree at $w$ (hence, especially, $t|_\eps = t$) and
 \item $t[t']_w$ is the tree that results from replacing $t|_w$ by $t'$.
\end{itemize}

\begin{definition}
 Let $\Sigma$ be an alphabet and $\square\notin\Sigma$. A \emph{1-context over
 $\Sigma$} is a tree $t\in U_\Sigma(\brc\square)$ such that $\square$ appears
 at exactly one position in $t$. The set of all such 1-contexts is denoted by
 $C_\Sigma$.
\end{definition}

1-contexts are used to describe trees that are not fully known yet: The
placeholder symbol $\sigma$ stands in for a missing subtree. Given a 1-context
$t\in C_\Sigma$ and a tree $t'\in U_\Sigma(V)$, we abbreviate $t[t'] :=
t[t']_w$ such that $t(w) = \square$. In other words, $t[t']$ is the tree that
results from $t$ when the $\square$ node is replaced by $t'$.

\begin{definition}
 A \emph{ranked alphabet} is a pair $(R,\rk)$ where $R$ is an alphabet and
 $\rk:R\to\zn$ is a mapping. $\rk$ is said to assign a \emph{rank} or
 \emph{arity} or each symbol in $R$.
\end{definition}

A ranked alphabet is usually denoted only by $R$. The existence of a suitable
$\rk$ is implied. A symbol $\rho\in R$ may be denoted as $\rho^{(k)}$ where
$\rk(\rho) = k$, to inform the reader of the symbol's rank.

\begin{definition}
 Let $R$ be a ranked alphabet and $V$ be a set. The \emph{set $T_R(V)$ of
 ranked trees over $R$ indexed by $V$} is defined by
 \[
  T_R(V) := \mbig\brc{t \in U_R(V) \dmiddle| \forall w\in\pos(t): \rk\mbig\kla{t(w)} = \text{number of children of $w$ in $t$}}.
 \]
\end{definition}

In other words, the number of children of each node in $t$ must be equal to the
rank of its label. For nodes labeled with symbols from $V$, a rank of 0 is
implied.

\begin{definition}
 Let $\Sigma$ be an alphabet. A \emph{regular tree grammar (RTG) over $\Sigma$} is a
 triple $\ug = (Q,q_0,R)$ where
 \begin{itemize}\setlength\itemsep{-0.3em}
  \item $Q$ is a nonempty alphabet (of states),
  \item $q_0\in Q$ is an initial state and
  \item $R\subset Q^*\times\Sigma\times Q$ is a finite ranked alphabet (of
   rules) such that every rule $\mbig\kla{(q_1,\ldots,q_k),\sigma,q}\in R$ has
   rank $k$.
 \end{itemize}
\end{definition}

We will write $\mbig\kla{(q_1,\ldots,q_k),\sigma,q}\in R$ as $q\to\sigma(q_1,\ldots,q_k)$ instead.

\begin{definition}
 Let $\ug=(Q,q_0,R)$ be an RTG over $\Sigma$. The \emph{family
 $\mbig\kla{D^q(\ug)\dmiddle|q\in Q}$ of partial abstract syntax trees oof
 $\ug$} is the smallest $Q$-indexed family $(D^q|q\in Q)$ such that, for all
 $q\in Q$,
 \[
  q \in D^q \quad\wedge\quad \forall \rho=\mbig\kla{q\to\sigma(q_1,\ldots,q_k)}\in R,d_i\in D^{q_i}: \rho(d_1,\ldots,d_k) \in D^q.
 \]
\end{definition}


An RTG describes a language (i.~e.~a countable set) of trees. Trees are derived
by starting with a tree containing only the a root node with the label $q_0$,
then successively replacing $Q$-labeled nodes according to the RTG's rule set
until no more $Q$-labeled nodes are left. The resulting tree is in $U_\Sigma$
and its abstract syntax tree is in $T_R$. Partial abstract syntax trees are in $T_R(Q)$.

For example, consider the RTG $\ug = \mbig\kla{\brc{q_0,q_1},q_0,R}$ where
$\Sigma=\brc{a,b,c}$ and
\[
 R = \brc{ q_0 \to a(q_0,q_1), \quad q_0 \to b, \quad q_1 \to c }.
\]
The tree $a(b,c)$ is derived like this: (Each row shows the partial tree $\in
U_\Sigma(Q)$ on the left and the partial abstract syntax tree $\in T_R(Q)$ on the
right.)
\[\begin{matrix}
 \begin{tikzpicture}[every node/.style={rectangle,draw,minimum size=2em},baseline=(current bounding box.center)]
  \node (n) at (0,0) { $q_0$ };
 \end{tikzpicture}
 &&
 \begin{tikzpicture}[every node/.style={rectangle,draw,minimum size=2em},baseline=(current bounding box.center)]
  \node (n) at (0,0) { $q_0$ };
 \end{tikzpicture}
 \\\\[-1em]
 \downarrow & \text{apply } q_0\to a(q_0,q_1) & \downarrow \\[0.5em]
 \begin{tikzpicture}[every node/.style={rectangle,draw,minimum size=2em},baseline=(current bounding box.center)]
  \node (n) at (0,0) { $a$ };
  \node (n1) at (-1,-1.5) { $q_0$ } edge (n);
  \node (n2) at (+1,-1.5) { $q_1$ } edge (n);
 \end{tikzpicture}
 &\hspace*{4em}&
 \begin{tikzpicture}[every node/.style={rectangle,draw,minimum size=2em},baseline=(current bounding box.center)]
  \node (n) at (0,0) { $q_0 \to a(q_0,q_1)$ };
  \node (n1) at (-1,-1.5) { $q_0$ } edge (n);
  \node (n2) at (+1,-1.5) { $q_1$ } edge (n);
 \end{tikzpicture}
 \\\\[-1em]
 \downarrow & \text{apply } q_0\to b & \downarrow \\[0.5em]
 \begin{tikzpicture}[every node/.style={rectangle,draw,minimum size=2em},baseline=(current bounding box.center)]
  \node (n) at (0,0) { $a$ };
  \node (n1) at (-1,-1.5) { $b$ } edge (n);
  \node (n2) at (+1,-1.5) { $q_1$ } edge (n);
 \end{tikzpicture}
 &\hspace*{4em}&
 \begin{tikzpicture}[every node/.style={rectangle,draw,minimum size=2em},baseline=(current bounding box.center)]
  \node (n) at (0,0) { $q_0 \to a(q_0,q_1)$ };
  \node (n1) at (-1,-1.5) { $q_0\to b$ } edge (n);
  \node (n2) at (+1,-1.5) { $q_1$ } edge (n);
 \end{tikzpicture}
 \\\\[-1em]
 \downarrow & \text{apply } q_1\to c & \downarrow \\[0.5em]
 \begin{tikzpicture}[every node/.style={rectangle,draw,minimum size=2em},baseline=(current bounding box.center)]
  \node (n) at (0,0) { $a$ };
  \node (n1) at (-1,-1.5) { $b$ } edge (n);
  \node (n2) at (+1,-1.5) { $c$ } edge (n);
 \end{tikzpicture}
 &\hspace*{4em}&
 \begin{tikzpicture}[every node/.style={rectangle,draw,minimum size=2em},baseline=(current bounding box.center)]
  \node (n) at (0,0) { $q_0 \to a(q_0,q_1)$ };
  \node (n1) at (-1,-1.5) { $q_0\to b$ } edge (n);
  \node (n2) at (+1,-1.5) { $q_1\to c$ } edge (n);
 \end{tikzpicture}
\end{matrix}\]

From the abstract syntax tree on the right, the tree on the left can be derived easily
by substituting the labels $q = \sigma(q_1,\ldots,q_k)\in R$ by the contained
symbols $\sigma\in\Sigma$. We shall call this projection $\pi_\Sigma: T_R(Q)\to U_\Sigma(Q)$.

\begin{definition}
 Let $\ug$ be an RTG over $\Sigma$. The \emph{language of $\ug$} is the set
 \[
  \lang\ug := \brc{t\in U_\Sigma | \exists d\in D^{q_0}(\ug): t=\pi_\Sigma(d)}.
 \]
 A language $L\subseteq U_\Sigma$ is \emph{recognizable} if there exists an RTG
 $\ug$ over $\Sigma$ such that $\lang\ug = L$.
\end{definition}

A useful property of RTGs is \emph{determinism}: An RTG $\ug$ over $\Sigma$ is
called \emph{deterministic} if, for any $(q_1,\ldots,q_k)\in Q^*$ and
$\sigma\in\Sigma$, there is at most one $q\in Q$ such that $\mbig\kla{q \to
\sigma(q_1,\ldots,q_k)} \in R$. A language is called deterministic if it is
described by a deterministic RTG.

From determinism results unambiguity: An RTG $\ug$ over $\Sigma$ is called
\emph{unambiguous} if, for every tree $t\in\lang\ug$, there exist exactly one
abstract syntax tree $d\in D^{q_0}(\ug)$ such that $\pi_\Sigma(d) = t$. To see why, have
a look at the fully derived tree $t = a(b,c)$ above. Since $\ug$ in this
example is deterministic, the abstract syntax tree can be recovered from $t$ by
traversing the nodes from the bottom up and assigning the rules that are used
at that position. At each position, we know the symbol $\sigma$ at this
position in the tree and the states $q_1,\ldots,q_k$ from the left sides of the
rules used for the child nodes. Therefore, the state $q$ (and therefore the
rule $\rho$) for this position can be chosen deterministically.

\begin{definition}
 A \emph{probabilistic regular tree grammar (PRTG) over $\Sigma$} is a pair
 $(\ug,p)$ of an RTG $\ug=(Q,q_0,R)$ over $\Sigma$ and a mapping
 $p:Q\to\zr_{\geq0}^{Q^*\times\Sigma}$ that is constrained to $R$ in the
 following way:
 \[
  \forall q\in Q, u\in Q^*\times\Sigma: p(q)(u) \neq 0 \Rightarrow (u,q) \in R.
 \]
 $(\ug,p)$ is called \emph{proper} if $p\in\um_R(Q^*\times\Sigma|Q)$. The
 \emph{meaning} of $(\ug,p)$ is the mapping
 \[
  \mbig\lang{(\ug,p)}: U_\Sigma \to \zr_0, \quad
  t \mapsto \sum_{d\in D^{q_0}(\ug): \pi_\Sigma(d)=t} p(d).
 \]
 Herein, $p(d) := p\mbig\kla{\pi(d)}$, where $\pi(d)$ is an $R$-corpus with
 \[
  \pi(d)(p) := \mbig\abs{\mbig\brc{w\in\pos(d): d(w)=p}}.
 \]
 The set of all PRTG over $\Sigma$ shall be denoted as $\ur(\Sigma)$.
\end{definition}

We usually write a PRTG $(\ug,p)$ as just $\ug$ and imply the existence of $p$.
The meaning function assigns a probability to a tree from $\lang\ug$ by summing
the probability of all derivations resulting in that tree, where the
probability of a derivation is the product of the probability of each rule
occurring in it. This implies that rule applications are statistically
independent from each other.

When describing hidden information in terms of PRTG in the next section, the
notion of \emph{inside and outside weights} will be useful to judge, broadly
speaking, how much a certain state contributes to the derivations of a certain
observation.

\begin{definition}
 Let $\ug = (Q,q_0,R)$ be a PRTG over $\Sigma$. The \emph{inside weight} of a
 state $q\in Q$ is given by
 \[
  \beta(q) := \sum_{d\in D^q(\ug)\cap T_R} p(d).
 \]
\end{definition}

The sum goes over all complete abstract syntax trees rooted at $q$. The inside weight
describes the collective probability of all derivations starting at $q$. Inside
weights are usually calculated by noting that each $d\in D^q(\ug)\cap T_R$ must
have a rule of the form $q\to\ldots$ at its root. The remaining probabilities
can then be expressed as the inside weights of the states that emerge this rule
application, giving
\[
 \beta(q) = \sum_{q_1,\ldots,q_k,\sigma} p\mbig\kla{q\to\sigma(q_1,\ldots,q_k)} \cdot \beta(q_1) \cdots \beta(q_k).
\]

The set of all such equations is a non-linear equation system in the $\beta(q)$
for $q\in Q$. In some cases, this system can be solved intuitively by starting
with those equations where no $\beta(q_i)$ occurs on the right side, then
substituting the obtained value in the other equations until they are all
solved. This is not possible, however, if states derive other states in a
cyclic way. For example:
\[
 R = \brc{\begin{aligned}
  \rho_1 &= q_0 \to \sigma_1(q_1, q_1), \\
  \rho_2 &= q_1 \to \sigma_2(q_0), \\
  \rho_3 &= q_1 \to \sigma_3
 \end{aligned}} \quad\leadsto\quad
 \begin{aligned}
  \beta(q_0) &= p(\rho_1) \cdot \beta(q_1)^2 \\
  \beta(q_1) &= p(\rho_2) \cdot \beta(q_0) + p(\rho_3)
 \end{aligned}
\]

In the general case, \cite[pp.~6]{bucstuvog15} shows that $\beta$ is the least fixpoint of the mapping $F: (\zr_{\geq0}\cup\brc\infty)^Q \to (\zr_{\geq0}\cup\brc\infty)^Q$, given by
\[
 F(u)(q) := \sum_{q_1,\ldots,q_k,\sigma} p\mbig\kla{q\to\sigma(q_1,\ldots,q_k)} \cdot u(q_1) \cdots u(q_k).
\]
Therefore,
\[
 \beta = \lim_{n\to\infty} F^n(u_0) \quad\text{where}\quad u_0(q) := 0 \;\forall q\in Q
\]
can be approximated by performing as many iterations of $F$ as desired.

\begin{definition}
 Let $\ug = (Q,q_0,R)$ be a PRTG over $\Sigma$. The \emph{outside weight} of a state $q\in Q$ is given by
 \[
  \alpha(q) := \sum_{d\in D^{q_0}(\ug): \exists d'\in C_R: d = d'[q]} p(d).
 \]
\end{definition}

The sum goes over all partial abstract syntax trees, where only one unexpanded
state is left, and that state is $q$. The outside weight describes the
collective probability of all derivations that use $q$, but without considering
at or below $q$. Similar to what we did with inside weights, we can expand this
definition into
\[
 \alpha(q) = \delta_{q_0}^q + \sum_{q',q_1,\ldots,q_k,\sigma,m:q_m=q} \alpha(q') \cdot p\mbig\kla{q'\to\sigma(q_1,\ldots,q_k)} \cdot \prod_{l\neq m} \beta(q_l)
\]

In this formulation, the sum goes over all rules that derive $q$, that is, $q$
is among the states $q_1,\ldots,q_k$ on the right hand side of the rule, at
index $m$. The outside weight $\alpha(q)$ considers the outside weight
$\alpha(q')$ of the previous state and the inside weight of all states adjacent
to $q$, but not the inside weight of $q$ itself.

Finally, we need to account for $q=q_0$, in which case the trivial tree $d$ containing only a $q_0$-labeled root node must be considered. Since $p(d)=1$ for this tree, it is accounted for in the above formula through the Kronecker symbol
\[
 \delta_{q_0}^q := \begin{cases}
  1 & \text{if } q_0 = q, \\
  0 & \text{otherwise}.
 \end{cases}
\]

Assuming that the inside weights $\beta(q)$ have already been calculated, the
equations for $\alpha(q)$ form a linear equation system that can be solved
efficiently with the standard algorithms for linear equation systems.

\section{Inside-outside step mapping}

We can now resume the task of describing hidden information in terms of trees,
such that the countable events, as introduced by the counting information, are
labels in these trees' nodes. Similar to how a language model can be described
as a counting information for use with the simple counting step mapping, a
language model may also be described as an inside-outside information, for use
with the inside-outside step mapping.

Again, we assume the previous definitions of $X$, $Y$, $c$, $A$, $B$ and $C$
throughout this section.

\begin{definition}
 An \emph{inside-outside (IO) information} is a quadruple\footnote{The definitions
 in this chapter have previously diverged from \cite{bucstuvog15} by omitting
 the set $Z$, which is only required for translation models, and otherwise set
 to $Z=\brc\emptyset$ for language models. For IO informations, this leads to
 additional divergence from the definition in \cite{bucstuvog15}, where $\mu$
 is introduced as a quintuple $\mu=(q,\pi_1,\pi_2,K,H)$. Both definitions are
 congruent for language models by letting $\pi_1 := \pi$ and $\pi_2(y) :=
 \emptyset\in Z$.} $\mu = (q,\pi,K,H)$ such that
 \begin{itemize}\setlength\itemsep{-0.3em}
  \item $C\subseteq A\times B$ is a ranked alphabet such that $Y_{\not\bot}\subseteq T_C$ contains ranked trees over $C$,
  \item $q:\Omega\to\um_C(A|B)$ as for counting informations,
  \item $\pi:Y_{\not\bot}\to X_{\not\bot}$ maps hidden information to observations,
  \item $K\in\ur(C)$ is an PRTG with $\lang K = Y_{\not\bot}$ and
  \item $H:X_{\not\bot}\to\ur(C)$ assigns a PRTG to every observation such that $\lang{H(x)} = \pi^{-1}(x)$.
 \end{itemize}
\end{definition}

Each hidden information $y$ is a tree from $T_C$, and a conditional probability
$q_\omega(a|b)$ can be assigned to each countable event $c = (a,b)$ that occurs
in the tree. Each hidden information $y$ belongs to exactly one observation
$x$. The set of all hidden informations $Y_{\not\bot}$ is described by the PRTG
$K$, and for each observation $x$, the set of all hidden informations leading
to this observation is described by the PRTG $H(x)$.

The description of a language model using an IO information is strictly more
powerful than the description via counting informations, since an IO
information induces a counting information.

\begin{definition}
 Let $\mu=(q,\pi,K,H)$ be an IO information. The \emph{induced counting
 information} $\mu^\flat = (q,\lambda,\pi)$ is given by
 \begin{align*}
  \lambda(x,y) &:= \begin{cases}
   1 & \text{if } \pi_1(y) = x, \\
   0 & \text{otherwise},
  \end{cases} \\
  \pi(x,y)(a,b) &:= \begin{cases}
   \mbig\abs{\mbig\brc{w\in\pos(y): y(w)=(a,b)}} & \text{if } \pi_1(y) = x, \\
   0 & \text{otherwise},
  \end{cases}
 \end{align*}
\end{definition}
The induced counting information $\mu^\flat$ is always proper.
\cite[p.~15]{bucstuvog15} Using it, we can construct a suitable complete-data corpus:
\[
 c\!\dangle{\omega,\mu}: C\to\zr_{\geq0}, \quad
 (a,b) \mapsto \sum_x c(x) \cdot \chi_{\omega,x}(a,b).
\]
Unlike $c\!\dangle{\omega,\varkappa}$, this definition uses the original corpus
$c$ instead of relying on the complete-data corpus of the previous step
mapping. In this expression, $(\chi_{\omega,x}|\omega\in\Omega,x\in
X_{\not\bot})$ is a family of $C$-corpora, defined by
\[
 \chi_{\omega,x}(a,b) := \beta(q_0)^{-1} \cdot \sum_{\rho = (q\to(a,b)(q_1,\ldots,q_k))\in R} \alpha(q) \cdot p(\rho) \cdot \beta(q_1) \cdots \beta(q_k),
\]
where $(Q,q_0,R)$ is the PRTG $H(x)$ with rule probabilities $p$ given by
\[
 p(\rho) := q_\omega(a|b) \quad\text{for}\quad \rho = \mbig\kla{q\to(a,b)(q_1,\ldots,q_k)}.
\]
To understand the formula for $\chi_{\omega,x}(a,b)$, recall that $H(x)$ is an
RTG such that its derivations, that is: the set of its complete abstract syntax
trees, is $\pi^{-1}(x)$, the set of all hidden informations that lead to the
observation $x$. Each of these derivations $d$ has a probability $p(d)$, and
the sum of all these probabilities is the inside weight $\beta(q_0)$.
$\xi_{\omega,x}(a,b)$ describes, in simple terms, the contribution of rules
containing the countable event $(a,b)$ to this total weight (such that
multiple occurrences of $(a,b)$ in a single derivation are counted multiple
times). For every such derivation, we have the part that leads up to the state
$q$ (described by the outside weight $\alpha(q)$), the rule $\rho$ that
produces $(a,b)$ with probability $p(\rho)$, and the derivations that occur
below this rule application (described by the inside weights $\beta(q_i)$ of
the produced states).

Given the complete-data corpus, we can once more apply a maximum-likelihood
estimator for $q$ to arrive at the \emph{inside-outside step mapping}:
\[
 \stepmap\mu_\mathrm{io}: \Omega_0 \to \up(\Omega), \quad
 \omega \mapsto \cmle_q\mbig\kla{c\mnorm\dangle{\omega,\mu}} = \argmax_{\omega'} q_{\omega'}\mbig\kla{c\mnorm\dangle{\omega,\mu}}.
\]
\cite[pp.~16]{bucstuvog15} shows that $\stepmap\mu_\mathrm{io} =
\mnorm\stepmap{\mu^\flat}_\mathrm{sc}$ for all IO informations $\mu$, from
which follows that $\stepmap\cdot_\mathrm{io}$ is nondecreasing.

\section{Review}

This chapter introduced the three different formalizations of language models
proposed by \cite{bucstuvog15}:
\begin{enumerate}
 \item models $p:\Omega\to\um(Y\times X)$
 \item counting informations $\varkappa=(q,\lambda,\pi)$
 \item inside-outside informations $\mu=(q,\pi,K,H)$
\end{enumerate}
Each formalization is stronger than the one before it, as evidenced by the
$\flat$ operator that converts each formalization into the previous one.
A step mapping can be defined for each of these formalizations.
\begin{center}\begin{tikzpicture}[scale=2,>=stealth]
 \begin{scope}[every node/.style={rectangle,draw},minimum height=2em]
  \node (sm) at (+0,+0) { step mapping };
  \node (cb) at (-2,+1) { model };
  \node (sc) at (+0,+1) { counting info };
  \node (io) at (+2,+1) { IO info };
 \end{scope}
 \begin{scope}[every node/.style={inner sep=0.5ex}]
  \draw[thick] (cb.south) edge[->,bend right=30] node [auto,swap,inner sep=0] {$\stepmap\cdot_\mathrm{cb}$} (sm.west);
  \draw[thick] (sc.south) edge[->] node [auto] {$\stepmap\cdot_\mathrm{sc}$} (sm.north);
  \draw[thick] (io.south) edge[->,bend left=30] node [auto] {$\stepmap\cdot_\mathrm{io}$} (sm.east);
  \draw[thick] (sc.west) edge[->] node [auto] {$\flat$} (cb.east);
  \draw[thick] (io.west) edge[->] node [auto] {$\flat$} (sc.east);
 \end{scope}
\end{tikzpicture}\end{center}

In the definition of both the counting information $\varkappa=(q,\lambda,\pi)$
and of the IO information $\mu=(q,\pi,K,H)$, the only part that depends on
$\omega$ is the probability model $q$. All other parts are fixed as soon as the
language model is instantiated, and are therefore not subject to training. For
many language models, it is sufficient to choose $\Omega := \um_C(A|B)$ and $q
= \operatorname{id}$ (i.~e.~$q_\omega = \omega$ for all $\omega$).

Applying lemma~\ref{lemma:empirical2}, we see that the empirical probability
distribution for the complete-data corpus is an element of
$\stepmap\mu_\mathrm{io}(\omega)$ (and similar for the simple counting step
mapping). Since the EM algorithm only needs one element from
$\stepmap\mu_\mathrm{io}(\omega)$, we can use
\[
 \stepmap\mu'_\mathrm{io}(\omega) := \mbig\brc{ \widetilde{c\mnorm\dangle{\omega,\mu}} }
\]
instead, which is much less expensive to compute than the full $\argmax$.

\chapter{The Hidden Markov model}

The definitions in this chapter are based on \cite{vogler2015} and largely
agree with most other definitions for Hidden Markov models in contemporary
literature, e.~g.~\cite[pp.~210]{jm09}, except for variable names.

\begin{definition}
 A \emph{Hidden Markov model} (HMM) is a quintuple $H = (Q,V,\#,t,e)$ such that
 \begin{itemize}\setlength\itemsep{-0.3em}
  \item $Q$ is a non-empty alphabet (of states),
  \item $V$ is a non-empty alphabet (of letters),
  \item $\#\notin Q\cup V$ is a separate (initial and final) state,
  \item $t\in\um(Q\cup\brc{\#}|Q\cup\brc{\#})$ and $e\in\um(V|Q)$.
 \end{itemize}
\end{definition}

{\color{red}TODO}


\appendix\addchap{Appendix}
\renewcommand\thesection{\Alph{section}}
\renewcommand\theequation{\Alph{section}.\arabic{equation}}
\renewcommand\thefigure{\Alph{section}.\arabic{figure}}
\renewcommand\thetable{\Alph{section}.\arabic{table}}
\setcounter{theorem}0
\renewcommand\thetheorem{\Alph{section}.\arabic{theorem}}

\section{Elided proofs from Chapter~3}\label{sec:appendix-hk-proofs}

This appendix contains the proofs for Lemmas~\ref{lemma:03-consistent-yield}
and~\ref{lemma:03-language-partition}. We imply that $\uh = (Q,V,\#,t,e)$ is an
HMM. Moreover, we imply all the definitions introduced in
Section~\ref{sect:03-deriving} to describe $\uh$ in terms of the IO information $\mu_\uh$.
As a prerequisite for both proofs, we shall begin by writing down the languages
of all the grammar states of $H(x)$ and $K$.

\begin{definition}
 Let $\ug = (Q_\ug,q_0,R_\ug)$ be an RTG over $C$. For any grammar state $q\in
 Q_\ug$, the \emph{language of $q$ generated by $\ug$} is given by
 \[
  \lang\ug_q := \pi_C\mbig\kla{D^q(\ug)\cap T_{R_\ug}}.
  \qedhere
 \]
\end{definition}

The mapping $\pi_C\colon T_{R_\ug}(Q_\ug)\to U_C(Q_\ug)$ was introduced as $\pi_\Sigma$ on
page~\pageref{def:02-pi-sigma}. The language $\lang\ug_q$ contains all trees
from $U_C$ that can be generated by applying the grammar's rules, starting from
the state $q$. It especially holds that
\begin{equation}\label{eq:a1}
 \lang\ug = \lang\ug_{q_0}.
\end{equation}
To find $D^q(\ug)\cap T_{R_\ug}$, we first observe that, according to
Definition~\ref{def:02-past}, each $d\in D^q(\ug)$ must have one of two forms,
\[
 d = q \quad\text{or}\quad d = \rho(d_1,\ldots,d_k),
\]
where $\rho=\mbig\kla{q\to c(q_1,\ldots,q_k)}\isin R_\ug$ such that $\rk(c) = k$, and
each $d_i\in D^{q_i}(\ug)$. The first form, $d = q$, is not allowed in
$\lang\ug_q$ since $q\notin T_{R_\ug}$. The second form is admissible, but
because $d\in T_{R_\ug}$, the subtrees $d_i$ may not have positions labeled with
a state $q\in Q_\ug$ either. We therefore have
\[
 D^q(\ug)\cap T_{R_\ug} = \bigcup_{\rho=(q\to c(q_1,\ldots,q_k))\isin R_\ug}
 \mbig\brc{\rho(d_1,\ldots,d_k) \npipe d_1\in D^{q_1}(\ug)\cap T_{R_\ug},\ldots, d_k\in D^{q_k}(\ug)\cap T_{R_\ug}},
\]
from which follows
\begin{equation}
 \label{eq:a2}
 \lang\ug_q = \bigcup_{\rho=(q\to c(q_1,\ldots,q_k))\isin R_\ug}
 \mbig\brc{c(t_1,\ldots,t_k) \npipe t_1\in\lang\ug_{q_1},\ldots, t_k\in\lang\ug_{q_k}},
\end{equation}
when $\pi_C$ is applied on both sides. We can thus describe the languages
$\mbig\lang K$ and $\mbig\lang{H(x)}$ recursively, in terms of the languages
generated from their states.

Let $x\in V^*$. Application of \eqref{eq:a2} to all reachable states in $H(x)$ yields
\begin{align*}
 \mbig\lang{H(x)}_{(T\!,q,x')} &= \bigcup_{q'\in Q}\mbig\brc{\mbig\kla{q',(q,T)}(t_1,t_2) \!\pipe t_1\in\mbig\lang{H(x)}_{(E,q',x')}, t_2\in\mbig\lang{H(x)}_{(T\!,q',\cdr(x'))}}, \\
 \mbig\lang{H(x)}_{(T\!,q,\eps)} &= \mbig\brc{\mbig\kla{\#,(q,T)}} \\
 \intertext{for all $q\in Q_\#$ and $x'\in\suff(x)$, and}
 \mbig\lang{H(x)}_{(E,q,x')} &= \mbig\brc{\mbig\kla{\car(x'),(q,E)}}
\end{align*}
for all $q\in Q$ and $x'\in\suff(x)$. Inserting the last equation into the one above it, we can describe $\mbig\lang{H(x)} = \mbig\lang{H(x)}_{(T\!,\#,x)}$ recursively by
\begin{subequations}\label{eq:a3}\begin{align}
 \mbig\lang{H(x)}_{(T\!,q,x')} &= \bigcup_{q'\in Q}\mBig\brc{\mbig\kla{q',(q,T)}\mBig\kla{\mbig\kla{\car(x'),(q',E)},t} \!\pipe t\in\mbig\lang{H(x)}_{(T\!,q',\cdr(x'))}}, \label{eq:a3a} \\
 \mbig\lang{H(x)}_{(T\!,q,\eps)} &= \mbig\brc{\mbig\kla{\#,(q,T)}} \label{eq:a3b}
\end{align}\end{subequations}
for all $q\in Q_\#$. Analogous application of~\eqref{eq:a2} to all states in
$K$, followed by insertion of the expression for $\lang K_{(E,q)}$ into that of
$\lang K_{(T\!,q)}$, results in
\begin{equation}\label{eq:a4}
 \lang K_{(T\!,q)} = \mbig\brc{\mbig\kla{\#,(q,T)}} \cup \bigcup_{q'\in Q}\bigcup_{v\in V}\mBig\brc{\mbig\kla{q',(q,T)}\mBig\kla{\mbig\kla{v,(q',E)},t} \!\pipe t\in\lang K_{(T\!,q')}}
\end{equation}
for any $q\in Q_\#$. An important observation can be made from the structure of~\eqref{eq:a3}.
\begin{equation}\label{eq:a5}
 \forall x\in V^+, q\in Q\colon\;
 \mbig\lang{H(x)}_{(T,q,\cdr(x))} = \mbig\lang{H\mbig\kla{\cdr(x)}}_{(T,q,\cdr(x))}.
\end{equation}
This equivalence holds because $Q_{\cdr(x)}$ and $R_{\cdr(x)}$ are strict
subsets of $Q_x$ and $R_x$, respectively, and furthermore, because $R_x$ does
not contain any rules not in $R_{\cdr(x)}$ that expand grammar states from
$Q_{\cdr(x)}$. Therefore, we have
\[
 D^q\mbig\kla{H(x)} = D^q\mBig\kla{H\mbig\kla{\cdr(x)}} \quad\forall q\in Q_{\cdr(x)},
\]
from which follows~\eqref{eq:a5}.

\subsection{Proof of Lemma~\ref{lemma:03-consistent-yield}}

\begin{repeatlemma}{\ref{lemma:03-consistent-yield}}
 With the definitions from Section~\ref{sect:03-deriving},
 \[
  \forall x\in X_{\not\bot}\colon\;
  \forall y\in \mbig\lang{H(x)}\colon\;
  \pi_1(y) = x.
 \]
\end{repeatlemma}
\clearpage

\emph{Proof.} Let $x\in X_{\not\bot} = V^*$. We are going to show that
\begin{equation}\label{eq:p11}
 \forall x'\in\suff(x)\cup\brc\eps\colon\;
 \forall q\in Q_\#\colon\;
 \pi_X\mbig\kla{\mbig\lang{H(x)}_{(T,q,x')}} = \brc{x'}.
\end{equation}
by induction over the length of $x'$. The lemma follows from this statement because
\[
 \pi_1\mbig\kla{\mbig\lang{H(x)}}
 = \pi_X\mbig\kla{\mbig\lang{H(x)}}
 = \pi_X\mbig\kla{\mbig\lang{H(x)}_{(T,\#,x)}} \overset{\eqref{eq:p11}}= \brc x.
\]

For the base case, let $x' = \eps$. Then~\eqref{eq:p11} holds because of
\[
 \pi_X\mbig\kla{\mbig\lang{H(x)}_{(T,q,\eps)}}
 \overset{\eqref{eq:a3b}}=
 \mbig\brc{\pi_X\mbig\kla{\#,(q,T)}} = \brc\eps
\]
for any $q\in Q_\#$. For the induction step, let $x'\in\suff(x)$. Since
$\abs{\cdr(x')} < \abs{x'}$, it can be assumed that $\cdr(x')$
satisfies~\eqref{eq:p11} (\emph{induction hypothesis}). Application of $\pi_X$
to both sides of~\eqref{eq:a3a} yields
\begin{align*}
 \pi_X\mbig\kla{\mbig\lang{H(x)}_{(T\!,q,x')}}
 &= \bigcup_{q'\in Q}\mBig\brc{\pi_X\mbig\kla{\car(x'),(q',E)}\pi_X(t) \!\pipe t\in\mbig\lang{H(x)}_{(T\!,q',\cdr(x'))}} \\
 &= \mbig\brc{\car(x')} \cdot \bigcup_{q'\in Q}\mBig\brc{\pi_X(t) \!\pipe t\in\mbig\lang{H(x)}_{(T\!,q',\cdr(x'))}} \\
 &= \mbig\brc{\car(x')} \cdot \bigcup_{q'\in Q}\underbrace{\pi_X\mbig\kla{\mbig\lang{H(x)}_{(T\!,q',\cdr(x'))}}}_{=\brc{\cdr(x')}\text{ by induction hypothesis}} \\
 &= \mbig\brc{\car(x')\cdr(x')} = \mbig\brc{x'}. \tag*{\qed}
\end{align*}

\subsection{Proof of Lemma~\ref{lemma:03-language-partition}}

\begin{repeatlemma}{\ref{lemma:03-language-partition}}
 With the definitions from Section~\ref{sect:03-deriving},
 \[
  \mbig\lang K = \bigcup_{x\in X_{\not\bot}} \mbig\lang{H(x)}.
 \]
\end{repeatlemma}

\emph{Proof.} Let $\pi_Q\colon T_C\to Q_\#$ be the mapping defined by
\[
 \pi_Q(y) := q \text{ such that } y(\eps) = \mbig\kla{a,(q,b)}.
\]
That is, $\pi_Q$ extracts the state $q$ from the label $\mbig\kla{a,(q,b)}$
of the root position of any $y$. As a prerequisite for the actual proof, we are
going to show that
\begin{equation}\label{eq:p21}
 \forall y\in T_C\colon\;
 \mBig\kla{y\in\mbig\lang K_{(T\!,\pi_Q(y))}}
 \text{ iff }
 \mBig\kla{y\in\mbig\lang{H\mbig\kla{\pi_X(y)}}_{(T\!,\pi_Q(y),\pi_X(y))}},
\end{equation}
by structural induction on $y$. According to the definition of the ranked
alphabet $C$, each $y\in T_C$ can have one of three forms,
\begin{align*}
 y_1(v,q) &:= \mbig\kla{v,(q,E)} &&\text{ where } v\in V\text{ and } q\in Q, \\
 y_2(q) &:= \mbig\kla{\#,(q,T)} &&\text{ where } q\in Q_\#, \\
 y_3(q,q',y_e,y_t) &:= \mbig\kla{q',(q,T)}(y_e,y_t) &&\text{ where } q\in Q_\#, q'\in Q, \text{ and } y_e,y_t\in T_C.
\end{align*}
The first two forms are the base cases for the structural induction on $y$,
and the last form represents the inductive step.

\subsubsection*{Base cases}

For any $y=y_1(v,q)$, we have
$\pi_Q(y)=q$ and $\pi_X(y)=v$. We see that~\eqref{eq:p21} holds because
\[
 \mbig\kla{v,(q,E)} \notin \mbig\lang K_{(T\!,q)}
 \quad\text{ and }\quad
 \mbig\kla{v,(q,E)} \notin \mbig\lang{H(v)}_{(T\!,q,v)}.
\]
For any $y=y_2(q)$, we have $\pi_Q(y) = q$ and $\pi_X(y)=\eps$, and \eqref{eq:p21} is
satisfied because
\[
 \mbig\kla{\#,(q,T)} \in \mbig\lang K_{(T\!,q)}
 \quad\text{ and }\quad
 \mbig\kla{\#,(q,T)} \in \mbig\lang{H(\eps)}_{(T\!,q,\eps)}.
\]

\subsubsection*{Induction step: ``$\Rightarrow$'' direction}

Consider any
\[
 y = y_3(q,q',y_e,y_t) = \mbig\kla{q',(q,T)}(y_e,y_t).
\]
such that $y\in\lang K_{(T,q)}$, and let $x:=\pi_X(y)$. Note that $\pi_Q(y) =
q$.  We see from~\eqref{eq:a4} that
\[
 y_t \in \mbig\lang K_{(T\!,q')}
 \quad\text{ and }\quad
 y_e = \mbig\kla{v,(q',E)}
 \quad\text{ where }\quad
 q'\in Q \text{ and } v\in V.
\]
From this follows $x = \pi_X(y) = v\pi_X(y_t)$ and, therefore, $\car(x) = v$
and $\cdr(x)=\pi_X(y_t)$. Since $y_t$ is a substructure of $y$, it can be
assumed to satisfy~\eqref{eq:p21} (\emph{induction hypothesis}). Since
$y_t\in\mbig\lang K_{(T\!,q')}$, we can see from~\eqref{eq:a4} that only
$\pi_Q(y_t) = q'$ is possible. Therefore, by the induction hypothesis,
\[
 y_t
 \in \mbig\lang{H\mbig\kla{\pi_X(y_t)}}_{(T\!,q',\pi_X(y_t))}
 = \mbig\lang{H\mbig\kla{\cdr(x)}}_{(T\!,q',\cdr(x))}.
\]

Because of~\eqref{eq:a5}, this is equivalent to saying that $y_t \in \mbig\lang{H(x)}_{(T\!,q',\cdr(x))}$. We therefore have
\[
 y = \mbig\kla{q',(q,T)}\mBig\kla{\mbig\kla{\car(x),(q',E)},y_t}
\]
with $q\in Q_\#$, $q'\in Q$, and $y_t \in \mbig\lang{H(x)}_{(T\!,q',\cdr(x))}$,
from which follows
\[
 y \in \mbig\lang{H(x)}_{(T\!,q,x)}
\]
because of~\eqref{eq:a3a}.

\subsubsection*{Induction step: ``$\Leftarrow$'' direction}

Consider any $y = y_3(q,q',y_e,y_t)$ such that
$y\in\mbig\lang{H(x)}_{(T\!,\pi_Q(y),x)}$, wherein $x := \pi_X(y)$. Since
$\pi_Q(y) = q$, it follows from~\eqref{eq:a3} that
$y_e = \mbig\kla{\car(x),(q',E)}$ and
$y_t \in \mbig\lang{H(x)}_{(T\!,q',\cdr(x))}$.
Because of~\eqref{eq:a5}, the second statement is equivalent to
\[
 y_t \in \mbig\lang{H\mbig\kla{\cdr(x)}}_{(T\!,q',\cdr(x))}.
\]
From~\eqref{eq:a3}, we see that $\pi_Q(y_t)$ must be $q'$. Since $y_t$ is a
substructure of $y$, it can be assumed to satisfy~\eqref{eq:p21}, from which
follows that $y_t \in \mbig\lang K_{(T\!,q')}$. According to~\eqref{eq:a4}, it
therefore holds that $y\in\mbig\lang K_{(T\!,q)}$. The proof for~\eqref{eq:p21}
is thereby complete.

\subsubsection*{Conclusion of the proof}

To finally derive Lemma~\ref{lemma:03-language-partition}, we first note that
\begin{equation}\label{eq:p22}
 \pi_Q\mbig\kla{\mbig\lang K_{(T\!,\#)}} = \brc\#
 \quad\text{ and }\quad
 \pi_Q\mbig\kla{\mbig\lang{H(x)}_{(T\!,\#,x)}} = \brc\#
\end{equation}
for every $x\in X_{\not\bot}$, which follows from~\eqref{eq:a4} and~\eqref{eq:a3}, respectively.
We can thus expand $\lang K$ into
\begin{align*}
 \mbig\lang K
 = \mbig\lang K_{(T\!,\#)}
 &= \mbig\brc{y\in \pi_Q^{-1}(\#) \npipe y\in\mbig\lang K_{(T\!,\#)}} \\
 &= \mbig\brc{y\in \pi_Q^{-1}(\#) \npipe y\in\mbig\lang K_{(T\!,\pi_Q(y))}}.
\end{align*}
By using~\eqref{eq:p21}, we obtain
\begin{align*}
 \mbig\lang K
 &= \mbig\brc{y\in \pi_Q^{-1}(\#) \npipe y\in\mbig\lang{H\mbig\kla{\pi_X(y)}}_{(T\!,\pi_Q(y),\pi_X(y))}} \\
 &= \mbig\brc{y\in \pi_Q^{-1}(\#) \npipe y\in\mbig\lang{H\mbig\kla{\pi_X(y)}}_{(T\!,\#,\pi_X(y))}}.
\end{align*}

To obtain the desired structure, we split this set into one set per yield $x = \pi_X(y)$.
\begin{align*}
 \mbig\lang K
 &= \bigcup_{x\in X_{\not\bot}} \mbig\brc{y\in \pi_Q^{-1}(\#) \npipe y\in\mbig\lang{H\mbig\kla{\pi_X(y)}}_{(T\!,\#,\pi_X(y))}, x = \pi_X(y)} \\
 &= \bigcup_{x\in X_{\not\bot}} \mbig\brc{y\in \pi_Q^{-1}(\#) \npipe y\in\mbig\lang{H(x)}_{(T\!,\#,x)}, x = \pi_X(y)}.
\end{align*}

Because of~\eqref{eq:p22}, each $H(x)_{(T,\#,x)}$ is a subset of $\pi_Q^{-1}(\#)$. We can therefore simplify this to
\[
 \mbig\lang K
 = \bigcup_{x\in X_{\not\bot}} \mbig\brc{y\in \mbig\lang{H(x)}_{(T\!,\#,x)} \npipe x = \pi_X(y)}
 = \bigcup_{x\in X_{\not\bot}} \mbig\brc{y\in \mbig\lang{H(x)} \npipe x = \pi_X(y)},
\]
which is equivalent to Lemma~\ref{lemma:03-language-partition} because of Lemma~\ref{lemma:03-consistent-yield}.\qed

\endinput
which is equivalent to the desired statement,
\[
 \mbig\lang K = \bigcup_{x\in X_{\not\bot}} \mbig\lang{H(x)},
\]
because of Lemma~\ref{lemma:03-consistent-yield}.\qed


\backmatter
% \listoffigures % TODO: necessary?
% \listofalgorithms % TODO: necessary?
\bibliographystyle{alpha}
\bibliography{citations}

\clearpage
\thispagestyle{empty}

\minisec{Erklärung}\vspace*{1.5em}

Hiermit erkläre ich, dass ich diese Arbeit im Rahmen der Betreuung im Institut
für Theoretische Informatik ohne unzulässige Hilfe Dritter verfasst und alle
Quellen als solche gekennzeichnet habe.

\vspace*{15em}

Stefan Majewsky \par
Dresden, {\color{red}August 2017} % TODO: enter final date when known

\end{document}
